% Generated draft: requires mathematical and source review.
\documentclass[11pt]{article}
\usepackage[margin=25mm]{geometry}
\usepackage{fontspec}
\setmainfont{Latin Modern Roman}
\usepackage{amsmath,amssymb,mathtools,unicode-math}
\setmathfont{Latin Modern Math}
\usepackage{enumitem,needspace,xurl,hyperref}
\setlength{\parindent}{0pt}
\setlength{\parskip}{5pt}
\setlength{\emergencystretch}{3em}
\newcommand{\N}{\mathbb{N}}
\newcommand{\Z}{\mathbb{Z}}
\newcommand{\Q}{\mathbb{Q}}
\newcommand{\R}{\mathbb{R}}
\newcommand{\C}{\mathbb{C}}
\newcommand{\F}{\mathbb{F}}
\newcommand{\A}{\mathbb{A}}
\newcommand{\PP}{\mathbb{P}}
\newcommand{\E}{\mathbb{E}}
\newcommand{\eps}{\varepsilon}
\newcommand{\dd}{\,\mathrm{d}}
\begin{document}
\section*{UnsolvedMath ID: 30000628}
\textbf{Source problem code:} catalogue entry
\subsection{Definitions and mathematical statement}
Extend the classical sphere-through-points and hyperplane-through-points incidence problems to convex bodies: determine existence and uniqueness of spheres, balls, or hyperplanes simultaneously tangent to, or slicing prescribed volume fractions from, well-separated families of convex bodies.

In a standard Euclidean setting, let $K_1,\ldots,K_m\subset\R^d$ be convex bodies, meaning compact convex sets with nonempty interior. Write $\operatorname{vol}_d$ for Lebesgue volume. A hyperplane with an indicated side has the form
\[
 H(u,t)=\{x:\langle u,x\rangle=t\},\quad
 H^+(u,t)=\{x:\langle u,x\rangle\ge t\},\quad |u|=1.
\]
A prescribed-fraction version of one requested incidence problem would solve
\[
 \operatorname{vol}_d(K_i\cap H^+(u,t))
       =a_i\operatorname{vol}_d(K_i)\qquad(1\le i\le m)
\]
for specified admissible fractions $a_i$. A ball version would instead use $K_i\cap B(c,r)$ with the same kind of prescribed volume constraints. The associated sphere is $\partial B(c,r)$. These equations distinguish the variants but do not assert their solvability for arbitrary data.

For tangency, a hyperplane ordinarily supports each relevant body and meets it; spherical or ball tangency needs an internal/external or containment convention. The source does not state one. It also does not define well-separated, give $m$ relative to $d$, prescribe general position or strict convexity, specify the allowed fractions or tangency signs, or say whether uniqueness is for oriented hyperplanes, unoriented hyperplanes, balls, spheres, or a selected branch of solutions.

Accordingly, the record is a collection of existence-and-uniqueness research questions rather than one fully quantified theorem. The sphere/ball and hyperplane cases, tangency and volume slicing, and existence and uniqueness remain distinct. No particular definition of well-separated, number of bodies, or universal uniqueness claim is selected to fill the omissions.
\subsection{Short English statement}
The source asks for several convex-body analogues of point-incidence problems, using tangency or prescribed-volume cuts by spheres, balls, and hyperplanes. The separation definition, number of bodies, allowed data, and uniqueness conventions are missing.
\subsection{Sources}
\begin{itemize}
\item Local UnsolvedMath registry, permanent ID 30000628, code catalogue entry, ``Incidence Problems for Families of Convex Bodies.'' Proposed year 2006; proposer unspecified.
\item \emph{Konvexgeometrie} (2007), Oberwolfach Reports, pp.\ 3321--3396, DOI 10.4171/owr/2006/56; \url{https://doi.org/10.4171/owr/2006/56}. The displayed report year 2007 differs from the DOI and proposed year 2006. This is the sole supplied literature reference.
\item Statement status exact; the supplied verification calls the multipart sentence clear and internally complete and says it was retained byte-for-byte. The missing technical hypotheses and variant choices still require review.
\item Literature triage checked 2026-08-21: open. It generically calls the item an unresolved problem/conjecture/direction and reports no authoritative complete resolution in the supplied targeted check, without assessing separate variants. Created 2026-08-25; updated 2026-09-16. Current status and source scope require independent verification.
\item Registry curation by Ulam AI / UnsolvedMath Contributors, CC BY 4.0; the underlying report retains its supplied provenance. Unreviewed statement-only draft.
\end{itemize}
\end{document}
